\documentclass[aspectratio=169,12pt]{beamer}
\usepackage{fontspec}
\usepackage{xeCJK}
\usepackage{booktabs}
\usepackage{graphicx}
\graphicspath{{../outputs/}}

\usetheme{Madrid}
\usecolortheme{default}
\setmainfont{Times New Roman}
\setCJKmainfont{Microsoft JhengHei}
\setCJKsansfont{Microsoft JhengHei}
\setbeamertemplate{navigation symbols}{}

\title{資訊部內部 AI Agent 方案}
\subtitle{知識導航、需求受理、結案知識回收}
\author{AI發展部　黃丞暘}
\date{2026 年 9 月 29 日}

\begin{document}
% 第 1 頁：開頭
\begin{frame}
\titlepage
\end{frame}

% 第 2 頁：三個建議的 Agent（What）
\begin{frame}{三個建議的 Agent｜What}
\footnotesize
\begin{columns}[T,onlytextwidth]
\begin{column}{0.26\textwidth}
\begin{block}{1. 知識導航助手}
\textbf{做什麼：}查詢新人指引、SOP、申請流程、系統窗口及開發／維運文件。\par\medskip
\textbf{輸出：}操作說明；無資料時明確表示無法確認。
\end{block}
\end{column}
\begin{column}{0.36\textwidth}
\begin{block}{2. 需求受理與草稿助手}
\textbf{做什麼：}整理口語問題或申請需求，反問缺少的必要資訊。\par\medskip
\textbf{輸出：}供人工確認的工單／申請草稿，不自動送件。
\end{block}
\end{column}
\begin{column}{0.32\textwidth}
\begin{block}{3. 結案知識回收助手}
\textbf{做什麼：}從已確認的處理紀錄產生 FAQ、Error Case 或 SOP 修訂草稿。\par\medskip
\textbf{輸出：}標示適用範圍與待確認處，主管審核。
\end{block}
\end{column}
\end{columns}
\vspace{0.5em}
\begin{alertblock}{範圍定位}
新人教育訓練是知識導航助手的第一個使用情境。
\end{alertblock}
\end{frame}

% 第 3 頁：主要原因（Why）
\begin{frame}{選擇這三項的原因｜Why}
\footnotesize
\begin{tabular}{p{0.20\textwidth}p{0.35\textwidth}p{0.37\textwidth}}
\toprule
\textbf{方案} & \textbf{目前問題} & \textbf{導入價值} \\
\midrule
知識導航 & 已有文件散落各處，新人與同仁仍需反覆詢問。 & 讓有效知識找得到、查得到來源與版本。 \\
\addlinespace[0.7em]
需求受理 & 使用者描述不完整，資訊部需反覆追問才能處理。 & 先補齊系統、環境、時間、影響範圍與證據。 \\
\addlinespace[0.7em]
知識回收 & 問題雖已解決，方法可能只留在個人經驗或對話中。 & 經審核後回流 FAQ／SOP，減少下次從頭處理。 \\
\bottomrule
\end{tabular}
\vspace{0.8em}
\begin{alertblock}{三者形成循環}
已核准知識供查詢 $\rightarrow$ 新需求完整受理 $\rightarrow$ 已解決經驗審核回收 $\rightarrow$ 更新知識。
\end{alertblock}
\end{frame}

% 第 4 頁：輔助說明（How）
\begin{frame}{測試與上線｜How}
\footnotesize
\begin{columns}[T,onlytextwidth]
\begin{column}{0.56\textwidth}
\begin{block}{測試}
\begin{enumerate}
\item \textbf{知識導航：}整理有效 SOP、負責人、版本及真實測試題；回答必須附來源。
\item \textbf{需求受理：}先收手動輸入，擷取欄位並反問缺漏，只產生草稿。
\item \textbf{知識回收：}手動提交少量已確認案例；整理成待審核 FAQ／SOP 草稿。
\item \textbf{迭代：}記錄錯誤、文件缺口及人工修正，再調整 Workflow。
\end{enumerate}
\end{block}
\end{column}
\begin{column}{0.41\textwidth}
\begin{alertblock}{正式環境上線}
\begin{itemize}
\item 僅使用公司核准的文件與資料。
\item 案例先去識別化。
\item 未確認的根因不得寫成 SOP。
\item 新知識須經主管核准才入知識庫。
\item 不自動寄信、送簽、派案、開權限或操作 PROD。
\end{itemize}
\end{alertblock}
\end{column}
\end{columns}
\end{frame}

% 第 5 頁：預期目標
\begin{frame}{預期目標與驗收方式}
\footnotesize
\begin{tabular}{p{0.22\textwidth}p{0.38\textwidth}p{0.32\textwidth}}
\toprule
\textbf{面向} & \textbf{量測方式} & \textbf{驗收條件} \\
\midrule
知識導航 & 可用答案取得時間、正確來源命中、有據回答與正確拒答。 & 扣除查證時間仍省時；答案可追溯。 \\
\addlinespace[0.6em]
需求受理 & 首次資訊完整率、補件輪數與草稿人工修正時間。 & 減少反覆詢問；送出前仍由人確認。 \\
\addlinespace[0.6em]
知識回收 & 草稿審核通過率、文件補齊時間、後續同類問題是否可查。 & 案例可複用；未驗證內容不入正式知識庫。 \\
\addlinespace[0.6em]
整體效益與風險 & 實際使用任務、覆核及維護工時、敏感資料事件。 & 有人使用，且風險可控。 \\
\bottomrule
\end{tabular}
\end{frame}

\end{document}
